\section{Discussion}

Separating embodied and AR interaction from the core ARC-ADHD paper improves causal interpretability. The core intervention can first establish whether the challenge-regulation-recovery mechanism is measurable and feasible. The present paper can then ask whether changing the action and environmental context alters the same mechanism.

The three-condition structure is important. A direct comparison of stationary training with AR would confound physical movement, camera use, spatial mapping, novelty, environmental search, response modality, and visual presentation. The E condition provides a necessary intermediate contrast. If A and E perform similarly while both differ from S, the parsimonious interpretation is an embodiment or movement effect rather than an AR-specific effect.

Novelty and engagement remain strong alternative explanations. AR may initially increase motivation without improving transferable regulation. Repeated exposure, expectancy or preference measurement, and independent transfer outcomes are therefore needed before assigning a mechanistic role to AR.

Home deployment introduces substantial heterogeneity. A successful laboratory AR task may fail in small rooms, poor lighting, cluttered environments, or on devices with different tracking characteristics. Under the programme's BYOD design, internal-validity protection instead depends on module-specific qualification, capability stratification, technical-failure rules, and sensitivity analyses. This increases implementation realism but makes hardware compatibility and exclusion part of the empirical feasibility question.

The embodied condition also requires careful construct language. Rotating or moving a phone is not automatically fine-motor training, and movement quantity is not a clinical benefit. Claims should be tied to measured sensorimotor control, cognitive non-inferiority, recovery, and independent transfer.

The most important negative result would be one in which AR is engaging and increases movement but does not improve recovery or transfer, or degrades cognitive control. Such a result would support retaining simpler embodied interactions while removing the AR layer.
